% ===== PRÉAMBULE CANONIQUE — à coller VERBATIM en tête de CHAQUE .tex =====
\documentclass[11pt,a4paper]{article}
\usepackage[utf8]{inputenc}\usepackage[T1]{fontenc}\usepackage{lmodern}
\usepackage[french]{babel}
\usepackage{amsmath,amssymb,mathtools}\usepackage{icomma}
\usepackage[version=4]{mhchem}          % \ce{...} : équations & formules
\usepackage{siunitx}\sisetup{locale=FR} % \qty{}{}, \si{}, \ang{}
\usepackage{chemfig}                    % \chemfig{...} : structures développées
\usepackage{xcolor}\usepackage{colortbl}
\usepackage{tikz}\usepackage{pgfplots}\pgfplotsset{compat=1.18}
\usepackage[most]{tcolorbox}\usepackage{enumitem}\usepackage{array,booktabs}
\usepackage[a4paper,margin=2.2cm]{geometry}
\usetikzlibrary{arrows.meta,calc,angles,quotes,positioning,patterns,backgrounds,shapes.geometric}
\definecolor{primary}{RGB}{222,49,99}\definecolor{primarydark}{RGB}{150,22,55}
\definecolor{defcol}{RGB}{0,114,178}\definecolor{methcol}{RGB}{52,92,148}
\definecolor{excol}{RGB}{0,135,90}\definecolor{attncol}{RGB}{200,40,55}
\tcbset{boxbase/.style={enhanced,breakable,boxrule=0.9pt,arc=2.5pt,
  left=3mm,right=3mm,top=2.3mm,bottom=2.3mm,fonttitle=\bfseries,coltitle=white}}
% --- boîtes TUTORIEL (noms EXACTS, ne pas renommer) ---
\newtcolorbox{cle}[1][]{boxbase,colback=defcol!6,colframe=defcol,
  title={\textsc{À retenir}\ifx\relax#1\relax\relax\else\textmd{ : \itshape #1}\fi}}
\newtcolorbox{methode}[1][]{boxbase,colback=methcol!7,colframe=methcol,sharp corners,
  title={\textsc{Méthode}\ifx\relax#1\relax\relax\else\textmd{ --- \itshape #1}\fi}}
\newtcolorbox{exemplebox}[1][]{boxbase,colback=excol!5,colframe=excol,
  title={\textsc{Exemple}\ifx\relax#1\relax\relax\else\textmd{ : \itshape #1}\fi}}
\newtcolorbox{piege}[1][]{boxbase,colback=attncol!6,colframe=attncol,
  title={\textsc{Piège}\ifx\relax#1\relax\relax\else\textmd{ : \itshape #1}\fi}}
\newtcolorbox{remarque}[1][]{boxbase,colback=black!4,colframe=black!45,
  title={\textsc{Remarque}\ifx\relax#1\relax\relax\else\textmd{ : \itshape #1}\fi}}
% --- boîtes EXERCICE / CORRIGÉ (noms EXACTS, auto-comptées) ---
\newtcolorbox[auto counter]{exo}[1][]{boxbase,colback=primary!4,colframe=primary,
  title={\textsc{Exercice}~\thetcbcounter\ifx\relax#1\relax\relax\else\textmd{ --- \itshape #1}\fi}}
\newtcolorbox[auto counter]{corrige}[1][]{boxbase,colback=excol!4,colframe=excol,
  title={\textsc{Corrigé}~\thetcbcounter\ifx\relax#1\relax\relax\else\textmd{ --- \itshape #1}\fi}}
\newcommand{\R}{\mathbb{R}}\newcommand{\N}{\mathbb{N}}\newcommand{\Z}{\mathbb{Z}}\newcommand{\ds}{\displaystyle}
% (un \usepackage{fancyhdr}+en-tête est possible ; il est ignoré côté web)
% =========================================================================
\begin{document}
% THEME_TITLE: Hydroxydes, sels et lecture des formules

\begin{cle}[un sel : un cation et un anion assemblés]
Un \textbf{sel} associe un \textbf{cation} (métal ou \ce{NH4+}) et un \textbf{anion} en une formule neutre.
On le nomme toujours \textbf{« [anion] de [cation] »} (l'anion d'abord). Cas particuliers :
\begin{itemize}[leftmargin=1.5em,topsep=2pt,itemsep=1pt]
  \item \textbf{hydroxyde} (une \emph{base}) : l'anion est \ce{OH-} $\Rightarrow$ « hydroxyde de [métal] » ;
  \item \textbf{sel acide} : il reste un \ce{H} $\Rightarrow$ préfixe « (di/mono)hydrogéno » sur l'anion ;
  \item \textbf{hydrure} : l'anion est \ce{H-} $\Rightarrow$ « hydrure de [métal] ».
\end{itemize}
\end{cle}

\begin{methode}[nommer / écrire un sel, un hydroxyde, un sel acide]
\textbf{Écrire} : identifier cation et anion (+ valence), \emph{croiser les valences}, simplifier,
parenthéser un ion polyatomique d'indice $\geq 2$. \\
\textbf{Nommer} : énoncer l'anion, « de », puis le cation (avec sa valence si elle varie). Pour un
\textbf{sel acide}, préfixer l'anion de « hydrogéno » (ou « dihydrogéno ») selon le nombre de \ce{H} restants :
\ce{NaHCO3} $=$ hydrogénocarbonate de sodium ; \ce{Na2HPO4} $=$ monohydrogénophosphate de sodium.
\end{methode}

\begin{cle}[lire les chiffres d'une formule]
Il faut distinguer trois sortes de chiffres :
\begin{itemize}[leftmargin=1.5em,topsep=2pt,itemsep=1pt]
  \item l'\textbf{indice} (petit, dans la formule) : nombre d'atomes d'une sorte \emph{dans un motif}
        (\ce{Na2CO3} : 2 \ce{Na}, 1 \ce{C}, 3 \ce{O}) ;
  \item le \textbf{coefficient} (grand, \emph{devant}) : multiplie \emph{toute} la formule
        (\ce{2 Na2CO3} $=$ 2 motifs) ;
  \item l'\textbf{eau de cristallisation} \ce{. $n$ H2O} (\emph{hydrate}) : $n$ molécules d'eau par motif
        (\ce{CuSO4 . 5 H2O}).
\end{itemize}
\end{cle}

\begin{exemplebox}[six lectures types]
\ce{NaCl} : chlorure de sodium ; \ce{K2SO4} : sulfate de potassium ; \ce{Ca(OH)2} : hydroxyde de
calcium ; \ce{Fe(NO3)3} : nitrate de fer(III) ; \ce{NaHCO3} : hydrogénocarbonate de sodium ;
\ce{CuSO4 . 5 H2O} : sulfate de cuivre(II) pentahydraté.
\end{exemplebox}

\begin{piege}[erreurs fréquentes au concours]
\begin{itemize}[leftmargin=1.4em,topsep=2pt,itemsep=1pt]
  \item \textbf{coefficient $\neq$ indice} : \ce{2 Na2CO3} représente \emph{2 molécules} (pas « 4 »), soit
        $2\times6=12$ atomes.
  \item l'écriture \ce{. $n$ H2O} n'ajoute \emph{pas} d'atomes au motif de base : \ce{CaCO3 . 2 H2O} reste
        un motif \ce{CaCO3} accompagné de 2 \ce{H2O}.
  \item dans le nom, l'\textbf{anion vient d'abord} : \ce{K2SO4} est « sulfate \emph{de} potassium ».
\end{itemize}
\end{piege}

\begin{remarque}[le vieux préfixe « bi- »]
On rencontre encore « bicarbonate » pour \ce{HCO3-} (hydrogénocarbonate) et « bisulfate » pour \ce{HSO4-}
(hydrogénosulfate) ; « bichromate de potassium » \ce{K2Cr2O7} désigne le \emph{dichromate}. La forme
« hydrogéno\ldots » (ou « di\ldots ») est préférable.
\end{remarque}

\end{document}
